\MajorSection{القسم الخامس}

\Couplet{١١٧}{لا خير ولا شر؛ إنهما أهواء الإرادة الفانية؛}{ما ينفعني أسمّيه خيرًا، وما يؤذيني ويضرني أعدّه شرًا.}{}

\Couplet{١١٨}{يتبدلان بتبدل المكان، ويتغيران بتغير الشعوب؛ وفي أقصر فسحة من الزمان؛}{ارتدت كل رذيلة تاج فضيلة، وحُظر كل خير بوصفه خطيئة أو جريمة.}{}

\Couplet{١١٩}{يتقاطعان ويتشابكان كخيوط ملتفة مختلطة، ويتصل هذا بذاك ويمتزج به؛}{ولا ترى إلا عين الخضر أين يبدأ أحدهما وأين ينتهي الآخر.}{\BurtonNote{حاشية برتون: يُفترض أن الخضر هو النبي إيليا.}}

\Couplet{١٢٠}{أي فانٍ يستطيع مصاحبة الخضر، وقد استدار موسى مذعورًا ليفر؛}{وأي إنسان يبصر مسبقًا الزهرة أو الثمرة، والقدر يجبره على غرس الشجرة؟}{}

\Couplet{١٢١}{واقرأ بدلًا من حرية إرادة الإنسان: الناموس الخالد، أنانكي، والقسمة، والمصير؛}{ما كان، وما هو كائن، وما سيكون أبدًا؛ النجم والحظ والقدر وأورد ونورن أو الضرورة.}{}

\Couplet{١٢٢}{«حالة الإنسان الطبيعية هي تدبير الله»؛ ذلك موضوع الحكيم السخيف؛}{«عصر الإنسان الأول كان عصرًا ذهبيًا»؛ ذلك حلم الشاعر في اليقظة.}{}

\Couplet{١٢٣}{وهم وجهل! قبل زمن طويل من أول نفَس تنفّسه الإنسان على الأرض؛}{كان العالم مشهدًا متصلًا من الألم والتعذيب والافتراس والموت.}{}

\Couplet{١٢٤}{حيث كانت الثدييات المرعبة في البرية تمزق أشباهها عضوًا عضوًا؛}{وحيث كانت الزواحف الهائلة في البحر تسبح في أمواج من الدم.}{}

\Couplet{١٢٥}{لم تكن «الأرض الفتية الجميلة» صالحة إلا لتفرخ سلالتها من الوحوش المخيفة؛}{تارة ملتهبة الحرارة، وتارة جامدة من الجليد، وتارة مبللة تندى بفيض بخاري.}{}

\Couplet{١٢٦}{وتلك الشمس البهية، النيّر الأكبر، «العريس» في قيثارة الملك؛}{منجم ملتهب يغلي وينفجر، كرة سوداء عابسة من نار تدور.}{}

\Couplet{١٢٧}{وذلك القمر اللطيف، النيّر الأصغر، مصباح العاشق وبهجة الراعي؛}{عالم مخرّب، وكرة احترقت وانطفأت، وجثة على طريق الليل.}{}

\Couplet{١٢٨}{قل، ماذا كان يكترث بالخير أو الشر ذلك الذي اتخذ جحره في ثقب الجبل؛}{ذلك الوحش الضاري المتغذي بالدم، الأشد وحشية من أشد الذئاب أو الدببة وحشية؟}{}

\Couplet{١٢٩}{وكم طال في أيام الإنسان السابقة لآدم أن يكون الأكل والشراب والنوم والتناسل}{كل حياة البهيمة ذات القدمين، حياة كاملة بلا مدوّنة أحكام ولا عقيدة؟}{}

\Couplet{١٣٠}{أفخر لباسه جلد كثيف الشعر، وأفخر أدواته شظية حجر؛}{وأجمل زينته جلد موشوم وثقوب يعلّق فيها قطعه من العظم.}{}

\Couplet{١٣١}{وكان يقاتل من أجل الأنثى كما يقاتل من أجل الطعام، حين توقظ مواسم أيار الرغبة الحارة؛}{وكانت تلك هي الشهوة التي نمت فصارت حبًا، حين أعارها الخيال نارًا أصفى.}{}

\Couplet{١٣٢}{فأين إذن «ناموس الطبيعة الأبدي الذي حفره الله على قلب الإنسان»؛}{تأمل رأسه القردي، وأقرّ أن هذا الشيء لم يكن يستطيع أداء دور أعلى.}{}

\Couplet{١٣٣}{ومع ذلك، مع توالي العصور الطويلة، تعلّم من القندس والقرد والنملة أن يبني}{مأوى للأب والأم والنسل، من الريح العاصفة واللهيب المؤذيين القاتلين.}{}

\Couplet{١٣٤}{وأخيرًا جاءت النار؛ حين ألقيت قطعة حجر في اللهب الذي يضيء جحره؛}{فأخرجت المعدن اللامع، وجعلت من سيد الوحوش سيدًا على البشر.}{}

\Couplet{١٣٥}{أما «الحس الأخلاقي»، عبارتك الزاهدية، فليس إلا عطية السنوات الأحدث؛}{ولد الضمير حين كان الإنسان قد تخلص من فرائه وذيله وأذنيه المدببتين.}{}

\Couplet{١٣٦}{أي ضمير عند الموري القاتل، الذي يقتل ضيفه بضربة غادرة؛}{سوى أسفه لأنه لا يستطيع قتل المزيد؟ وأي وخزة توبة يمكن أن يعرف؟}{}

\Couplet{١٣٧}{تقولون إن «قسوة الأشياء» لغز لعينكم الكليلة؛}{وهي، إذ تثبت على نقطة في الفضاء، تغفل عن المشروع العام.}{}

\Couplet{١٣٨}{فانظروا! مضى الماموث في طريقه وصار ذكرى واسمًا؛}{أما ذو اليد، الذي يعقل نصف عقل، فقد بقي ليطالب برتبته ومكانه.}{\BurtonNote{حاشية برتون: الفيل.}}

\Couplet{١٣٩}{الزلزال والطاعون، والعاصفة والقتال والاشتباك، لا بد أن يعدّها الإنسان نذرًا ولعنات؛}{إذ لا ينظر إلا إلى نفسه، ولا يكترث بتتبّع المدى والخطة.}{}

\Couplet{١٤٠}{والزلزال الذي يأتي في طرفة عين ليخرّب ويسوّي بالأرض ويبتلع ويقتل؛}{يبني عالمًا لاستعمال أفضل، ويجعل الشر الخاص في خدمة الخير العام.}{}

\Couplet{١٤١}{وأفظع صوت يمكن أن تسمعه أذن الإنسان، حرب الريح العاصفة واندفاعها؛}{ينقّي مادة الحياة الإنسانية، ويولّد الصحة والقوة لبني البشر.}{}

\Couplet{١٤٢}{فماذا تسمونها: خيرات أم شرورًا، خيرات شريرة أم شرورًا خيّرة، خسارة أم ربحًا؛}{حين تقوم ممالك وينهار سقف، ويُكتسب عالم ويُقتل إنسان؟}{}

\Couplet{١٤٣}{وهكذا يجري شوط الوجود، إلى أن تبدّل الأرض قطبها، ربما في زمن مقبل؛}{ويرى أهل المشتري نجمًا آخر يسقط.}{\BurtonNote{حاشية برتون: كوكب المشتري.}}

\Couplet{١٤٤}{يرونه يسقط ويغيب عن البصر؛ من أين جاء وإلى أين ذهب، لا يستطيع الفكر أن يخبر؛}{اشرب من ذلك الجدول السرابي، وطارد رنين جرس الجمل!}{}
